\section{Syntax}

Let $\mathit{Tag} \ni t$ be a countably infinite set of \emph{tags}.

Outputs are asynchronous, sums are guarded, and recursion is as in VACCS, with
process variables $X$ and a tagged recursion operator:
\[
\begin{array}{r@{\ ::=\ }l}
P & \tagged{t}{c!e} \mid K \mid P \parallel P \mid \ite{e}{P}{P} \mid P \setminus c
    \mid X \mid \tagged{t}{\rec{X}{P}}\\
K & \mathbf{0}_{\tilde e} \mid \alpha.P \mid K + K\\
\alpha & \tagged{t}{\tau} \mid \tagged{t}{c?x} \mid
        \tagged{t}{U(\tilde e)} \mid \tagged{t}{\M(\tilde e \rhd x)}
        \mid \tagged{t}{\mathsf{init}(\tilde e)}
\end{array}
\]
The prefix $\tagged{t}{U(\tilde e)}$ applies the quantum gate $U$ to the
qubits $\tilde e$, and the prefix $\tagged{t}{\mathsf{init}(\tilde e)}$ resets
the qubits $\tilde e$ to $|0\rangle$. Both are super-operators
(\cref{def:sop}), and the memory treats them alike.
$P[Q/X]$ is the substitution of $Q$ for the free occurrences of $X$ in $P$;
unfolding a recursion is a silent move, selected by the tag of the recursion.
A \emph{choice} is either a tag or an unordered pair of tags,
$s ::= t \mid \{t,t'\}$: the pair $\{t,t'\}$ selects the synchronisation
between a move tagged $t$ and a move tagged $t'$, whichever side each move is
on. Throughout, $\lambda$ ranges over all
actions, $\tau$ included, and $\mu, \mu_1, \mu_2, \ldots$ over external actions.

The LTS of lqCCS configurations $\conf{\rho}{P}$ is given in \cref{fig:conf}.


\begin{figure}[h]
\small
\begin{mathpar}
\inferrule*[left=Tau]
  { }
  {\conf{\rho}{\tagged{t}{\tau}.P} \step{\tau}{t} \conf{\rho}{P}}

\inferrule*[left=Op]
  {U : \mathsf{Op}(|\tilde q|)}
  {\conf{\rho}{\tagged{t}{U(\tilde q)}.P} \step{\tau}{t} \conf{\upd{U}}{P}}

\inferrule*[left=Meas]
  {\M : \mathsf{Meas}(|\tilde q|) \\ 0 \le m < |\M|}
  {\conf{\rho}{\tagged{t}{\M(\tilde q \rhd x)}.P} \step{\tau}{t} \conf{\br}{P[m/x]}}

\inferrule*[left=Init]
  { }
  {\conf{\rho}{\tagged{t}{\initp}.P} \step{\tau}{t} \conf{\inirho}{P}}

\inferrule*[left=Send]
  {e \Downarrow v}
  {\conf{\rho}{\tagged{t}{c!e}} \step{c!v}{t} \conf{\rho}{\mathbf{0}}}

\inferrule*[left=Recv]
  { }
  {\conf{\rho}{\tagged{t}{c?x}.P} \step{c?v}{t} \conf{\rho}{P[v/x]}}

\inferrule*[left=IteT]
  {e \Downarrow \mathit{tt} \\ \conf{\rho}{P} \step{\lambda}{s} \conf{\rho'}{P'}}
  {\conf{\rho}{\ite{e}{P}{Q}} \step{\lambda}{s} \conf{\rho'}{P'}}

\inferrule*[left=IteF]
  {e \Downarrow \mathit{ff} \\ \conf{\rho}{Q} \step{\lambda}{s} \conf{\rho'}{Q'}}
  {\conf{\rho}{\ite{e}{P}{Q}} \step{\lambda}{s} \conf{\rho'}{Q'}}

\inferrule*[left=SumL]
  {\conf{\rho}{K_1} \step{\lambda}{s} \conf{\rho'}{P'}}
  {\conf{\rho}{K_1 + K_2} \step{\lambda}{s} \conf{\rho'}{P'}}

\inferrule*[left=SumR]
  {\conf{\rho}{K_2} \step{\lambda}{s} \conf{\rho'}{P'}}
  {\conf{\rho}{K_1 + K_2} \step{\lambda}{s} \conf{\rho'}{P'}}

\inferrule*[left=ParL]
  {\conf{\rho}{P} \step{\lambda}{s} \conf{\rho'}{P'}}
  {\conf{\rho}{P \parallel Q} \step{\lambda}{s} \conf{\rho'}{P' \parallel Q}}

\inferrule*[left=ParR]
  {\conf{\rho}{Q} \step{\lambda}{s} \conf{\rho'}{Q'}}
  {\conf{\rho}{P \parallel Q} \step{\lambda}{s} \conf{\rho'}{P \parallel Q'}}

\inferrule*[left=SynchL]
  {\conf{\rho}{P} \step{c!v}{t} \conf{\rho}{P'} \\
   \conf{\rho}{Q} \step{c?v}{t'} \conf{\rho}{Q'}}
  {\conf{\rho}{P \parallel Q} \step{\tau}{\{t,t'\}} \conf{\rho}{P' \parallel Q'}}

\inferrule*[left=SynchR]
  {\conf{\rho}{P} \step{c?v}{t} \conf{\rho}{P'} \\
   \conf{\rho}{Q} \step{c!v}{t'} \conf{\rho}{Q'}}
  {\conf{\rho}{P \parallel Q} \step{\tau}{\{t,t'\}} \conf{\rho}{P' \parallel Q'}}

\inferrule*[left=Restr]
  {\conf{\rho}{P} \step{\lambda}{s} \conf{\rho'}{P'} \\
   \forall v.\ \lambda \neq c!v, c?v}
  {\conf{\rho}{P \setminus c} \step{\lambda}{s} \conf{\rho'}{P' \setminus c}}

\inferrule*[left=Rec]
  { }
  {\conf{\rho}{\tagged{t}{\rec{X}{P}}} \step{\tau}{t}
   \conf{\rho}{P[\tagged{t}{\rec{X}{P}}/X]}}
\end{mathpar}
\caption{Configuration LTS, $\lambda \in \{\tau, c!v, c?v\}$. A measurement
moves to one unnormalised branch for each outcome $m$; all branches carry the
same choice $t$.}
\label{fig:conf}
\end{figure}


\begin{figure}[h]
\small
\begin{mathpar}
\inferrule*[left=Tau]
  { }
  {\tagged{t}{\tau}.P \step{\tau}{t} P}

\inferrule*[left=Op]
  {U : \mathsf{Op}(|\tilde q|)}
  {\tagged{t}{U(\tilde q)}.P \step{\req{U}}{t} P}

\inferrule*[left=Meas]
  {\M : \mathsf{Meas}(|\tilde q|) \\ m \in \mathbb{N}}
  {\tagged{t}{\M(\tilde q \rhd x)}.P \step{\mrecv{m}}{t} P[m/x]}

\inferrule*[left=Init]
  { }
  {\tagged{t}{\initp}.P \step{\initp}{t} P}

\inferrule*[left=Send]
  {e \Downarrow v}
  {\tagged{t}{c!e} \step{c!v}{t} \mathbf{0}}

\inferrule*[left=Recv]
  { }
  {\tagged{t}{c?x}.P \step{c?v}{t} P[v/x]}

\inferrule*[left=IteT]
  {e \Downarrow \mathit{tt} \\ P \step{\lambda}{s} P'}
  {\ite{e}{P}{Q} \step{\lambda}{s} P'}

\inferrule*[left=IteF]
  {e \Downarrow \mathit{ff} \\ Q \step{\lambda}{s} Q'}
  {\ite{e}{P}{Q} \step{\lambda}{s} Q'}

\inferrule*[left=SumL]
  {K_1 \step{\lambda}{s} P'}
  {K_1 + K_2 \step{\lambda}{s} P'}

\inferrule*[left=SumR]
  {K_2 \step{\lambda}{s} P'}
  {K_1 + K_2 \step{\lambda}{s} P'}

\inferrule*[left=ParL]
  {P \step{\lambda}{s} P'}
  {P \parallel Q \step{\lambda}{s} P' \parallel Q}

\inferrule*[left=ParR]
  {Q \step{\lambda}{s} Q'}
  {P \parallel Q \step{\lambda}{s} P \parallel Q'}

\inferrule*[left=SynchL]
  {P \step{c!v}{t} P' \\ Q \step{c?v}{t'} Q'}
  {P \parallel Q \step{\tau}{\{t,t'\}} P' \parallel Q'}

\inferrule*[left=SynchR]
  {P \step{c?v}{t} P' \\ Q \step{c!v}{t'} Q'}
  {P \parallel Q \step{\tau}{\{t,t'\}} P' \parallel Q'}

\inferrule*[left=Restr]
  {P \step{\lambda}{s} P' \\ \forall v.\ \lambda \neq c!v, c?v}
  {P \setminus c \step{\lambda}{s} P' \setminus c}

\inferrule*[left=Rec]
  { }
  {\tagged{t}{\rec{X}{P}} \step{\tau}{t} P[\tagged{t}{\rec{X}{P}}/X]}
\end{mathpar}
\caption{Program LTS, $\lambda \in \{\tau\} \cup \{c!v, c?v, \req{U}, \mrecv{m}, \initp\}$.
Quantum operations are requests to the memory: $\req{U}$ asks for $U$ on
$\tilde q$, $\mrecv{m}$ receives the outcome $m$ of $\M$ on $\tilde q$,
$\initp$ asks for the initialisation of $\tilde q$. The first and the last
are both served by the action $\memop{\sop{E}}$ of the memory.}
\label{fig:prog}
\end{figure}

\begin{figure}[h]
\small
\begin{mathpar}
\inferrule*[left=Sop]
  {\sop{E} : \mathsf{Sop}(|\tilde q|)}
  {\rho \mto{\memop{\sop{E}}} {\sapp{\sop{E}}}}

\inferrule*[left=Meas]
  {\M : \mathsf{Meas}(|\tilde q|) \\ 0 \le m < |\M|}
  {\rho \mto{\msend{m}} {\br}}
\end{mathpar}
\caption{Memory LTS, $\mu \in \{\memop{\sop{E}}, \msend{m}\}$. The memory has
two kinds of actions only: it applies a super-operator, or it sends the outcome
of a measurement.}
\label{fig:mem}
\end{figure}
